Fix \(0<e<1\), a treatment arm \(a\), and a standard Borel outcome space \(\mathcal Y\).   First omit the mutual-absolute-continuity restriction and let \(Q\in\mathfrak M^{\to}_e(P)\).  By consistency,
\[
 \mathcal L(Y(a)\mid A=a)=\mathcal L(Y\mid A=a)=P.
\]
Define the probability law \(R=\mathcal L(Y(a)\mid A\ne a)\).  The law of total probability and \(\Pr(A=a)=e\) give
\[
 Q=\mathcal L(Y(a))=eP+(1-e)R. \tag{1}
\]
Thus \(Q\in\mathfrak L^{\to}_e(P)\).  Notice that no support condition on \(R\) was used.

Conversely, suppose \(Q=eP+(1-e)R\) for an arbitrary probability law \(R\) on \(\mathcal Y\).  Since \(\mathcal Y\) is standard Borel, one may construct a joint probability law as follows.  Draw \(A\) with \(\Pr(A=a)=e\); conditionally on \(A=a\), draw \(Y(a)\) with law \(P\); conditionally on \(A\ne a\), draw \(Y(a)\) with law \(R\).  Draw \(Y(1-a)\), conditionally independently given \(A\), from any fixed probability law on \(\mathcal Y\).  Put
\[
 U=(A,Y(0),Y(1)),
\]
let the treatment structural equation return the first coordinate of \(U\), and let the outcome structural equation be \(Y=Y(A)\).  The common latent input \(U\) represents \(A\leftrightarrow Y\), and the explicit treatment input in the outcome equation represents \(A\to Y\).  Hence the construction is Markov to the bow ADMG in \(\mathrm{ass:bow\text{-}structure}\).  It satisfies consistency, has propensity \(e\), has observed arm law \(P\), and (1) gives interventional law \(Q\).  Therefore \(Q\in\mathfrak M^{\to}_e(P)\), proving
\[
 \mathfrak M^{\to}_e(P)=\mathfrak L^{\to}_e(P). \tag{2}
\]

We next prove the cap characterization.  If (1) holds, then \(Q(B)\ge eP(B)\) for every measurable \(B\).  Consequently \(P\ll Q\).  For \(r=dP/dQ\), the measure inequality \(eP\le Q\) implies \(er\le1\), \(Q\)-almost surely; otherwise integrating \(er-1\) on the positive-measure set where it is positive would contradict \(eP\le Q\).  Thus \(r\le1/e\).  Conversely, suppose \(P\ll Q\) and \(r=dP/dQ\le1/e\), \(Q\)-almost surely.  Define
\[
 R(B)=\frac{Q(B)-eP(B)}{1-e}.
\]
The numerator has nonnegative \(Q\)-density \(1-er\), so \(R\) is a nonnegative countably additive measure.  Moreover,
\[
 R(\mathcal Y)=\frac{1-e}{1-e}=1,
\]
and hence \(R\) is a probability law and \(Q=eP+(1-e)R\).  Together with (2), this proves
\[
 \mathfrak M^{\to}_e(P)=\mathfrak L^{\to}_e(P)
 =\left\{Q:P\ll Q,\ \frac{dP}{dQ}\le\frac1e\quad Q\text{-almost surely}\right\}. \tag{3}
\]